\section{符号说明}
容量与搬运量均使用题面容量单位；时间及工作量使用时钟周期。节点编号、序列位置和地址偏移为整数。缓存类型集合记为 $\mathcal T$。本文用 $F_v$ 表示结束时间，以免与边集 $E$ 混淆。
{\small\begin{longtable}{ll}\caption{主要符号及取值}\label{tab:symbols}\\\toprule
符号 & 定义与取值\\\midrule\endfirsthead
\multicolumn{2}{c}{表\thetable\ （续）}\\\toprule 符号 & 定义与取值\\\midrule\endhead\bottomrule\endfoot
$G=(V,E),\ V_o$ & 原始有向无环图、非管理操作节点集合\\
$N, B, B_v$ & 原始节点数、缓冲区集合、操作 $v$ 引用的去重缓冲区集合\\
$s_b,t_b,C_t$ & 正整数缓冲区大小、类型及该类型容量\\
$a_b,f_b$ & 缓冲区 $b$ 唯一的原始申请、释放节点\\
$\pi(v)$ & 节点位置；问题一为 $V$ 到 $\{1,\ldots,N\}$ 的双射\\
$x_{bj},o_{bj},z_{bj}$ & 第 $j$ 个驻留段的非负整数偏移、开始与结束节点\\
$\chi_b$ & 被任一 COPY\_IN 引用时取 1，否则取 0\\
$\mathcal M, M$ & 换出事件集合及事件数；同一缓冲区可重复出现\\
$H,D,T$ & 最大逻辑驻留量、额外搬运量、总执行周期\\
$c_v,p_v$ & 非负执行周期、指定执行单元；管理节点无独占单元\\
$S_v,F_v,r_v$ & 非负开始时间、结束时间、下游最长路径长度\\
$L_H,L_0,L_R$ & 同时使用容量下界、原始工作下界、固定驻留图时间下界\\
$W,\alpha$ & 就绪窗口大小、启发式评分权重\\
\end{longtable}}
